\documentclass[12pt]{article}
\usepackage[utf8]{inputenc}
\usepackage[T1]{fontenc}
\usepackage[margin=1in]{geometry}
\usepackage{setspace}
\usepackage{xcolor}
\usepackage[colorlinks=true, urlcolor=blue!60!black, linkcolor=blue!60!black]{hyperref}
\setlength{\parindent}{0em}
\setlength{\parskip}{0.6em}
\renewcommand{\baselinestretch}{1.0}
\begin{document}
\begin{center}
{\LARGE \textit{Clinton v. City of New York}, 524 U.S. 417 (1998)\par}
\vspace{0.5em}
{\large Case Brief --- Two-Page Edition\par}
\end{center}

\section{Memory Jogger}
The Line Item Veto Act violates the Presentment Clause: canceling parts of signed statutes is a unilateral partial repeal without bicameral passage and presentment (\href{https://supreme.justia.com/cases/federal/us/524/417/}{524 U.S. 417}).

\section{Detailed Case Facts}
The 1996 Act let the President, within five days of signing a bill, cancel spending items and limited tax benefits, which then lost ``legal force or effect'' unless Congress passed a disapproval bill. President Clinton canceled a provision forgiving up to \$2.6 billion in New York Medicaid taxes and a capital-gains deferral for sales of processing plants to farmers' cooperatives.

\section{Procedural History}
New York City and hospital groups, and the Idaho cooperative Snake River Potato Growers, sued; the district court held the Act unconstitutional, and the Supreme Court affirmed on expedited appeal.

\section{Judicial Votes}
Affirmed, 6--3. Stevens wrote for Rehnquist, Kennedy, Souter, Thomas, and Ginsburg; Kennedy concurred; Scalia, with O'Connor, concurred in part and dissented in part; Breyer dissented.

\section{Holding}
Both sets of plaintiffs had standing, and the Act's cancellation procedures violate Article I, \S~7 because they let the President amend or repeal parts of enacted statutes; any change must come through Article V.

\section{Analysis of Opinions}
Stevens, following \textit{INS v. Chadha}, held that cancellation was a partial repeal. A constitutional veto comes before a bill becomes law and covers the whole bill; cancellation comes after and strikes a part, and constitutional silence is ``equivalent to an express prohibition.'' \textit{Field v. Clark} involved contingent, mandatory action executing Congress's policy. Kennedy stressed that separation of powers protects individual liberty. Scalia saw a delegation of the long-accepted power to ``decline to spend'' and quipped that the Act's title ``succeeded in faking out the Supreme Court.'' Breyer argued that canceling under a statutory proviso executes the law and satisfies the ``intelligible principle'' test.

\section{Examples: Future Applications}
\textbf{Same-Side:} a power to cancel earmarks within ten days of signing is invalid. \textbf{Opposite-Side:} expedited rescission votes that still require both Houses and the President's signature are valid. \textbf{Fence-Sitter:} directives barring spending under ``sums not exceeding'' appropriations, which the majority conceded are permissible.

\section{Critique}
Critics following Scalia ask whether the line between ``declining to spend'' and erasing statutory text elevates drafting over substance. The Court also left open severability and whether canceling tax benefits differs from declining to spend.

\section{Key Quotations}
\begin{quote}
``Failure of political will does not justify unconstitutional remedies\ldots\ Liberty is always at stake when one or more of the branches seek to transgress the separation of powers\ldots\ Abdication of responsibility is not part of the constitutional design.''
\end{quote}
\noindent --- Justice Kennedy, concurring, 524 U.S. at 449, 450, 452.

\end{document}
